\chapter{Conception}
\label{chap:conception}

\section{Introduction}
\label{sec:co-intro}

Ce chapitre a pour objectif de traduire, dans le contexte du projet, les enseignements de l'état de l'art en une conception structurée, argumentée et directement exploitable pour la phase de réalisation. Le diagnostic qui l'ouvre ne reprend donc pas l'analyse théorique précédente~: il en propose une lecture opérationnelle afin d'identifier les lacunes à traiter dans l'architecture cible. Le chapitre présente ensuite la spécification des besoins fonctionnels et non fonctionnels, puis une vue d'ensemble de la solution proposée. Il expose enfin les principaux artefacts de pilotage et de modélisation du projet~: la planification des sprints, le backlog produit, les diagrammes UML et la planification temporelle, qui constituent la référence de travail pour la mise en \oe{}uvre.

\section{Diagnostic de l'existant}
\label{sec:analyse-existant}

\subsection{État des lieux de la sécurité IoT dans les déploiements opérateur}

L'analyse conduite en amont de ce projet, combinant l'étude des pratiques observées dans des déploiements IoT typiques et les enseignements publiés par l'ENISA~\cite{enisa_iot} et le NIST~\cite{nist_sp800213}, met en évidence plusieurs lacunes de sécurité récurrentes dans ce type d'environnement~:

\begin{itemize}[leftmargin=*]
	\item \textbf{Transport non systématiquement chiffré}~: le protocole MQTT est fréquemment déployé sur le port 1883 (texte clair), sans TLS, par souci de simplicité ou de performance perçue. Un attaquant positionné sur le segment réseau peut alors capturer l'intégralité des messages publiés, y compris les données métier sensibles~;
	\item \textbf{Authentification unidirectionnelle ou absente}~: le broker ne vérifie généralement pas le certificat du client. Les dispositifs s'authentifient souvent par un simple couple identifiant/mot de passe, voire sans aucun mécanisme d'authentification, rendant toute usurpation d'identité triviale~;
	\item \textbf{Absence de PKI dédiée}~: lorsque des certificats existent, ils sont fréquemment auto-signés, non révocables automatiquement et gérés de façon ad hoc, sans cycle de vie défini ni procédure de révocation d'urgence en cas de compromission~;
	\item \textbf{Contrôle d'accès applicatif insuffisant}~: les listes de contrôle d'accès aux topics MQTT sont rarement mises en œuvre, permettant à tout dispositif connecté d'accéder à l'ensemble des canaux de publication et de souscription, indépendamment de son rôle légitime~;
	\item \textbf{Supervision fragmentée}~: les outils de monitoring réseau et applicatif ne communiquent pas entre eux, rendant difficile la corrélation d'événements de sécurité. La détection d'une attaque nécessite une analyse manuelle croisée de plusieurs sources de journaux hétérogènes~;
	\item \textbf{Absence de détection comportementale}~: les mécanismes de détection d'intrusion existants se limitent aux signatures connues. Ils ne permettent pas de détecter des anomalies comportementales subtiles telles que l'utilisation d'un certificat valide depuis un contexte inhabituel, une dérive de fréquence de publication ou une exfiltration lente via des topics légitimes.
\end{itemize}

\subsection{Synthèse du diagnostic}

Le tableau~\ref{tab:diagnostic} synthétise ce diagnostic en mettant en regard chaque lacune identifiée avec la brique de sécurité retenue dans la solution, établissant ainsi le lien direct entre le contexte opérationnel observé et les choix de conception du projet.

\begin{table}[!htbp]
\centering
\caption{Diagnostic des lacunes et briques de réponse correspondantes}
\label{tab:diagnostic}
\small
\begin{tabularx}{\linewidth}{|>{\raggedright\arraybackslash}X|>{\raggedright\arraybackslash}X|}
\hline
\textbf{Lacune identifiée} & \textbf{Brique de réponse retenue} \\
\hline
Flux MQTT en clair, sans chiffrement de bout en bout & Activation obligatoire de TLS~1.3 sur Mosquitto (port~8883) \\ \hline
Authentification client absente ou basée sur mot de passe faible & Certificat X.509 client obligatoire, vérifié par le broker en mTLS \\ \hline
Gestion manuelle et non automatisée des certificats & PKI HashiCorp Vault à deux niveaux avec émission et révocation via API \\ \hline
Accès applicatifs non restreints par dispositif & ACL Mosquitto fondées sur le \emph{Common Name} du certificat client \\ \hline
Supervision réseau et applicative déconnectées & Intégration de Suricata au SIEM Wazuh pour la corrélation \\ \hline
Absence de détection comportementale & Module ML fondé sur IsolationForest et Random Forest \\ \hline
\end{tabularx}
\end{table}

Ce diagnostic structuré constitue la justification directe des choix techniques opérés dans ce projet~: chaque composant de la solution répond à une lacune précisément identifiée. Cette correspondance garantit la pertinence de l'architecture proposée au-delà d'une simple accumulation d'outils technologiques.

\section{Spécification des besoins}
\label{sec:co-besoins}

\subsection{Besoins fonctionnels (BF)}

\begin{table}[!htbp]
\centering
\caption{Besoins fonctionnels}
\label{tab:bf}
\begin{tabularx}{\linewidth}{|c|>{\raggedright\arraybackslash}X|}
\hline
\textbf{ID} & \textbf{Description} \\
\hline
\acs{BF-1} & Émettre, renouveler et révoquer automatiquement des certificats X.509 pour services et dispositifs. \\ \hline
\acs{BF-2} & Imposer le chiffrement TLS~1.3 et l'authentification mutuelle (mTLS) sur le broker MQTT. \\ \hline
\acs{BF-3} & Restreindre l'accès aux topics MQTT par dispositif via des listes de contrôle d'accès. \\ \hline
\acs{BF-4} & Simuler une flotte hétérogène de 50 dispositifs IoT générant un trafic réaliste. \\ \hline
\acs{BF-5} & Détecter les attaques réseau et applicatives par règles IDS et corréler les alertes dans un SIEM. \\ \hline
\acs{BF-6} & Détecter les anomalies non couvertes par les règles via un modèle d'apprentissage automatique. \\ \hline
\end{tabularx}
\end{table}

\subsection{Besoins non fonctionnels (BNF)}

\begin{table}[!htbp]
\centering
\caption{Besoins non fonctionnels}
\label{tab:bnf}
\begin{tabularx}{\linewidth}{|c|>{\raggedright\arraybackslash}X|}
\hline
\textbf{ID} & \textbf{Description} \\
\hline
\acs{BNF-1} & \textbf{Sécurité}~: aucun flux IoT en clair . rotation des secrets . séparation CA racine / intermédiaire. \\ \hline
\acs{BNF-2} & \textbf{Performance}~: latence handshake mTLS \(<50\) ms . débit broker \(>500\) msg/s par client. \\ \hline
\acs{BNF-3} & \textbf{Reproductibilité}~: l'environnement complet doit être réinstallable à partir des scripts fournis. \\ \hline
\acs{BNF-4} & \textbf{Scalabilité}~: l'architecture doit accepter sans refonte au moins 200 dispositifs. \\ \hline
\acs{BNF-5} & \textbf{Maintenabilité}~: configurations versionnées, modules indépendants, journalisation centralisée. \\ \hline
\acs{BNF-6} & \textbf{Observabilité}~: tableau de bord Wazuh accessible aux analystes en moins de 3 clics. \\ \hline
\end{tabularx}
\end{table}

Ces besoins définissent le \emph{quoi} de la solution~: les fonctionnalités attendues et les contraintes opérationnelles à respecter. Le diagramme de cas d'utilisation est présenté en trois parties, une par acteur de la plateforme.

\subsection{Diagramme de cas d'utilisation}

Trois acteurs principaux interagissent avec la plateforme~: le dispositif IoT, l'administrateur de sécurité et l'opérateur SOC. Les relations d'inclusion signalent les sous-cas obligatoires, tandis que les relations d'extension représentent les comportements optionnels. La connexion MQTT impose mTLS et TLS~1.3. L'IDS et le module ML complètent ensuite la détection.

\paragraph{Partie 1 — Dispositif IoT.}
La figure~\ref{fig:uc-p1} résume quatre actions du dispositif IoT~: demande de certificat X.509, connexion mTLS au broker, publication de télémétrie et renouvellement automatique du certificat.

\begin{figure}[!htbp]
\centering
\includegraphics[width=0.88\linewidth]{diag-uc-part1.png}
\caption{Cas d'utilisation --- Partie~1~: Dispositif IoT (authentification, publication, renouvellement)}
\label{fig:uc-p1}
\end{figure}
\FloatBarrier

\paragraph{Partie 2 — Administrateur sécurité.}
La figure~\ref{fig:uc-p2} montre les quatre cas d'utilisation de l'administrateur~: configuration de la PKI Vault (hiérarchie CA, rôles, TTL), définition des règles ACL Mosquitto, paramétrage des règles IDS Suricata et déploiement des règles de corrélation SIEM Wazuh.

\begin{figure}[!htbp]
\centering
\includegraphics[width=0.72\linewidth]{diag-uc-part2.png}
\caption{Cas d'utilisation --- Partie~2~: Administrateur sécurité (PKI, ACL, IDS, SIEM)}
\label{fig:uc-p2}
\end{figure}
\FloatBarrier

\paragraph{Partie 3 — Opérateur SOC.}
La figure~\ref{fig:uc-p3} représente les quatre cas d'utilisation de l'opérateur~: consultation du tableau de bord Wazuh, analyse des alertes IDS, lancement de la détection ML d'anomalies et évaluation des performances de la plateforme.

\begin{figure}[!htbp]
\centering
\includegraphics[width=0.72\linewidth]{diag-uc-part3.png}
\caption{Cas d'utilisation --- Partie~3~: Opérateur TT/SOC (dashboard, alertes, ML, performances)}
\label{fig:uc-p3}
\end{figure}
\FloatBarrier

\section{Vue d'ensemble de la solution proposée}
\label{sec:solution-proposee}

Pour répondre à la problématique et atteindre les objectifs ci-dessus, nous proposons une \textbf{architecture de défense en profondeur} composée de six briques complémentaires, intégralement déployée dans un laboratoire simulé sous \acs{GNS3}~\cite{gns3_docs} et Docker~\cite{docker_docs}~:

\begin{enumerate}[leftmargin=*]
	\item Une \textbf{infrastructure réseau segmentée} (VLANs IoT, services, monitoring, management) construite autour d'un pare-feu pfSense et d'un routeur MikroTik~.
	\item Une \textbf{PKI automatisée} HashiCorp Vault à deux niveaux (CA racine hors ligne + CA intermédiaire en ligne), capable d'émettre à la demande des certificats X.509 pour les services et les dispositifs~\cite{vault_pki}~.
	\item Un \textbf{broker MQTT Mosquitto} configuré en \acs{TLS}~1.3 avec authentification mutuelle obligatoire et listes de contrôle d'accès par topic~\cite{mosquitto_tls}~.
	\item Un \textbf{simulateur de flotte} de 50 dispositifs Python rejouant 76\,000 messages issus d'un jeu de données réel, dont sept catégories d'attaques~.
	\item Une \textbf{chaîne de détection} comprenant un IDS Suricata avec règles MQTT personnalisées~\cite{suricata_mqtt} et un SIEM Wazuh~\cite{wazuh_docs} pour la corrélation et la visualisation~.
	\item Un \textbf{pipeline de \acs{ML}} (IsolationForest~\cite{liu2008isolation} et Random Forest~\cite{breiman2001random}) entraîné sur les caractéristiques extraites du trafic, capable de détecter les anomalies non couvertes par les règles statiques.
\end{enumerate}

Le tableau~\ref{tab:mapping-sol} établit la correspondance entre les lacunes diagnostiquées à la section~\ref{sec:analyse-existant}, les briques de la solution et les objectifs spécifiques auxquels elles répondent, assurant ainsi la traçabilité complète entre le diagnostic initial et l'architecture proposée.

\begin{table}[!htbp]
\centering
\caption{Correspondance lacunes -- briques de sécurité -- objectifs spécifiques}
\label{tab:mapping-sol}
\small
\begin{tabularx}{\linewidth}{|>{\raggedright\arraybackslash}X|l|c|}
\hline
\textbf{Brique de sécurité} & \textbf{Section de réalisation} & \textbf{Obj.} \\
\hline
Topologie réseau segmentée (GNS3, pfSense, MikroTik) & §\ref{sec:re-archi} & O1 \\ \hline
PKI Vault à deux niveaux (CA racine + CA intermédiaire) & §\ref{sec:re-pki} & O2 \\ \hline
Broker Mosquitto TLS~1.3 + mTLS + ACL par topic & §\ref{sec:re-mqtt} & O3 \\ \hline
Simulateur de flotte de 50 dispositifs hétérogènes & §\ref{sec:re-attaques} & O4 \\ \hline
IDS Suricata avec règles MQTT dédiées & §\ref{sec:re-attaques} & O5 \\ \hline
SIEM Wazuh (corrélation d'alertes et tableau de bord) & §\ref{sec:re-attaques} & O6 \\ \hline
Module ML (IsolationForest + Random Forest) & §\ref{sec:re-ml} & O7 \\ \hline
\end{tabularx}
\end{table}

La figure~\ref{fig:archi-vue-intro} schématise cette vision globale, qui sera détaillée et justifiée dans les chapitres suivants.

\begin{figure}[!htbp]
\centering
\includegraphics[width=\textwidth]{diag-architecture-generale.png}
\caption{Vue d'ensemble de l'architecture de sécurisation proposée}
\label{fig:archi-vue-intro}
\end{figure}

La démarche Scrum retenue au chapitre~\ref{chap:contexte-general} structure la réalisation en six sprints~; la section suivante en détaille la planification et les livrables.

\section{Planification des sprints}
\label{sec:co-sprints}

Six sprints de deux semaines structurent la phase de réalisation. Le tableau~\ref{tab:sprints} en présente les objectifs principaux, du cadrage réseau initial jusqu'à l'évaluation du module ML.

\begin{table}[!htbp]
\centering
\caption{Planification des sprints}
\label{tab:sprints}
\small
\begin{tabularx}{\linewidth}{|l|l|>{\raggedright\arraybackslash}X|}
\hline
\textbf{Sprint} & \textbf{Durée} & \textbf{Objectif principal} \\
\hline
S1 & 2~semaines & Cadrage réseau, topologie segmentée, conteneurs de base et plan d'adressage. \\ \hline
S2 & 2~semaines & Déploiement de Vault, CA racine et intermédiaire, rôles PKI, scripts d'émission. \\ \hline
S3 & 2~semaines & Configuration de Mosquitto en TLS~1.3+mTLS, ACL par certificat et tests d'authentification. \\ \hline
S4 & 2~semaines & Simulation de 50 clients Python, génération de trafic et rejeu de 76\,000 messages. \\ \hline
S5 & 2~semaines & Règles Suricata MQTT, intégration Wazuh et tableaux de bord. \\ \hline
S6 & 2~semaines & Modèles IsolationForest et RandomForest, évaluation finale et rédaction. \\ \hline
\end{tabularx}
\end{table}

Cette planification met en évidence une progression cohérente~: l'infrastructure et la PKI précèdent logiquement le broker sécurisé, puis la simulation, la détection et enfin l'évaluation ML.

\section{Backlog produit}
\label{sec:co-backlog}

Le \emph{backlog produit} a été priorisé selon la méthode \textbf{MoSCoW} puis réparti sur les six sprints. Les éléments jugés critiques portent sur la topologie réseau, la PKI Vault, le broker mTLS, la simulation de flotte, l'intégration Suricata/Wazuh et l'évaluation du module ML. Le détail exhaustif des \emph{User Stories} n'est pas reproduit ici afin de conserver le focus du chapitre sur les choix de conception.

\section{Modélisation UML}
\label{sec:co-uml}

La modélisation UML a été recentrée sur trois vues suffisantes pour comprendre la solution~: la structure du simulateur, l'émission d'un certificat client et le cycle de publication mTLS.

\subsection{Diagramme de classes du simulateur}

La figure~\ref{fig:classe-p1} résume le noyau d'orchestration du simulateur~: \texttt{FleetOrchestrator} pilote les instances \texttt{DeviceSimulator}, tandis que \texttt{DeviceStats} centralise les métriques utiles à l'évaluation.

\begin{figure}[!htbp]
\centering
\includegraphics[width=0.75\linewidth]{diag-classe-part1.png}
\caption{Diagramme de classes du simulateur (orchestration et collecte des métriques)}
\label{fig:classe-p1}
\end{figure}
\FloatBarrier

\subsection{Diagramme de séquence~: émission d'un certificat client}

La figure~\ref{fig:seq-pki-p2} résume l'émission d'un certificat client via Vault~: requête API authentifiée, contrôle du rôle PKI, signature par la CA intermédiaire puis retour du bundle au dispositif, qui peut ensuite programmer son renouvellement.

\begin{figure}[!htbp]
\centering
\includegraphics[width=0.88\linewidth]{diag-pki-seq-part2.png}
\caption{Séquence d'émission d'un certificat client via Vault (\textasciitilde 87~ms)}
\label{fig:seq-pki-p2}
\end{figure}
\FloatBarrier

\subsection{Diagramme d'activité~: cycle de publication mTLS}

La figure~\ref{fig:activity-p3} synthétise le comportement d'un dispositif simulé~: chargement de son sous-ensemble de données, publication MQTT sur les topics autorisés, adaptation du rythme d'émission en cas d'attaque, puis agrégation des statistiques en fin d'exécution.

\begin{figure}[!htbp]
\centering
\includegraphics[width=0.72\linewidth]{diag-activity-part3.png}
\caption{Diagramme d'activité du cycle de publication mTLS}
\label{fig:activity-p3}
\end{figure}
\FloatBarrier

\section{Conclusion}
\label{sec:co-conclusion}

Ce chapitre a formalisé la conception du projet en articulant successivement le diagnostic de l'existant, la spécification des besoins, la vue d'ensemble de la solution retenue, la planification générale et les principales vues UML. Cet ensemble fournit une base cohérente pour la phase de réalisation, en assurant la traçabilité entre les lacunes identifiées, les objectifs visés et les choix de conception adoptés. Le chapitre suivant décrit précisément \textbf{comment} chaque élément a été mis en \oe{}uvre~: de l'environnement GNS3 à la PKI Vault, du broker mTLS au pipeline ML, en passant par les scénarios d'attaque.
